\section{Introduction}\label{sec:introduction}

For a normed space $Y$, a vector $x\in Y$, and $R>0$, consider the
symmetric lens of displacements
\[
 \mathcal L(x,R)=\{y\in Y:\|x+y\|\le R,\ \|x-y\|\le R\}.
\]
When $\|x\|$ is close to $R$, convexity restricts this set. In an
infinite-dimensional space one can ask for smallness after projecting
away finitely many coordinates. We prove such an estimate for a
particular tree norm, with a bound determined by the squared-radius
deficit $R^2-\|x\|^2$ of the center.

Approximate-midpoint lenses are central to asymptotic midpoint uniform
convexity: Dilworth, Kutzarova, Randrianarivony, Revalski, and Zhivkov
characterized that property in terms of the measure of noncompactness
of suitable lenses
\cite[Theorem~2.1 and Definition~2.2]{DKRRZ2016}. Their framework concerns
finite-codimensional control of symmetric displacements. The estimate
here gives explicit control outside each finite ancestral set for the
specified norm. Its proof uses disjoint tree segments, following the
norm construction introduced by James \cite[p.~739]{James1974}.

\subsection{The space and the tail estimate}

For each integer $h\ge1$, let $T_h=\mathbb N^{\le h}$ be the
rooted tree of words of length at most $h$, ordered by extension of
words. Its root is the empty word. Keep separate copies of these trees
in the forest
\[
 \mathcal F=\coprod_{h\ge1}(\{h\}\times T_h).
\]
Write $a\preceq b$ if $a$ is an ancestor of $b$ in the same component.
A \emph{segment} is a nonempty interval
$[a,b]=\{v:a\preceq v\preceq b\}$ in an ancestral chain.
For a finitely supported real function $u$ on $\mathcal F$, set
$u(S)=\sum_{v\in S}u(v)$ and define
\begin{equation}\label{eq:segment-norm}
 \rho(u)=\sup\left\{\left(\sum_{i=1}^m u(S_i)^2\right)^{1/2}:
  m\ge1,\ S_1,\ldots,S_m\text{ pairwise disjoint segments}\right\}.
\end{equation}
The Euclidean triangle inequality and singleton segments show that
$\rho$ is a norm. Let $J$ be its completion and let $X=J^*$, with
the dual norm denoted by $\|\cdot\|$. We write $e_v\in J$ for a
coordinate unit vector. Singleton tests give $|u(v)|\le\rho(u)$,
so coordinate evaluation extends to an element $e_v^*\in X$.
The construction in \eqref{eq:segment-norm} uses James's
square-sum norm on disjoint segments \cite[p.~739]{James1974}; the
branching and component heights here differ from the classical binary
tree.


The second model uses the infinite tree $T_\infty=\mathbb N^{<\omega}$,
including its root. Let $J_\infty$ be the completion for the same
disjoint-segment norm, and let
\[
 B_\infty=\overline{\operatorname{span}}\{e_v^*:v\in T_\infty\}
 \subset J_\infty^*.
\]
This closure uses the dual norm. It is the coordinate predual in the
Ghoussoub--Maurey--Schachermayer construction presented by
Banakh \cite[Section~2.II]{Banakh2000}. Section~\ref{binf:definition}
proves its finite-coordinate and atomic descriptions separately.
Perreau asked whether the canonical predual of the countably branching
James tree space is asymptotically midpoint uniformly convex
\cite[Question~6, p.~23 of arXiv v1]{Perreau2021}. This question
motivates our study of the infinite-height coordinate predual $B_\infty$
with the real norm specified here; the finite-height forest is a
separate model.

A finite set $H$ is \emph{ancestral} if it contains every ancestor of
its vertices. On either coordinate space write
\[
 P_Hx=\sum_{v\in H}x(e_v)e_v^*,\qquad Q_Hx=x-P_Hx.
\]
Both projections are contractive. The displacement below may have a
nonzero head.

\begin{theorem}\label{thm:main-tail}\label{mol:prop:joint-budget}
Let $Y$ be either the finite-height forest dual $X=J^*$ or the
infinite-height coordinate predual $B_\infty$. For every finite
ancestral $H$, every $x,y\in Y$ with $x=P_Hx$, and every $R\ge0$
with $\|x+y\|,\|x-y\|\le R$,
\begin{equation}\label{eq:main-squared-tail}
 \|Q_Hy\|\le2\sqrt{R^2-\|x\|^2}.
\end{equation}
\end{theorem}

The same head works for every displacement in the lens. This is the
feature needed when a single center has infinitely many possible
midpoints. Corollary~\ref{mol:cor:all-radius-averaged} gives an
explicit averaged asymptotic midpoint modulus for both models.

For comparison with the tree-space history, Girardi proved asymptotic uniform convexity (AUC) of the
canonical predual and full dual of the classical binary James tree
space \cite[Theorems~3 and~5 of arXiv v1]{Girardi2001}. Our finite-height model uses
countably branching trees with finite, unbounded heights. The local
cancellation estimate in this paper is proved for those specified
segment tests; it is not imported from the binary-tree result.

\subsection{The mechanism of the proof}

Section~\ref{sec:forest} proves the coordinate structure locally, so
no companion theorem is needed to read the argument. The tail splits
into an $\ell_2$ sum over the descendant subtrees starting just outside
a finite ancestral set.

In Section~\ref{sec:atoms}, we represent points of each finite head
unit ball by convex combinations of disjoint segment functionals.
A functional nearly norming the center makes the atom coefficients close
to its sums over the same segments. The difficulty is cancellation:
two represented atoms can have opposite sums when they leave the head
at the same vertex. Their head segments are nested. Splitting the
longer one improves its squared segment test, so the opposite-sign
product is bounded by the amount that the two tests fall short of one.
Section~\ref{sec:cancellation} proves this crossing comparison.

The two endpoint distributions can be combined with a random sign.
Keeping the coefficient error and the unused head-test energy in one
budget gives Theorem~\ref{thm:main-tail}, proved in
Section~\ref{mol:sec:joint-budget}. The same head controls every
displacement, so finite-coordinate approximation yields an asymptotic
midpoint modulus uniformly over all centers.

The obstruction to an equivalent AUC norm uses a different feature of
the spaces. Every root-path functional has norm one, while its child
increments are unit coordinate vectors and are weakly null. A positive
one-sided asymptotic modulus would permit successive children whose
path norms grow by a fixed factor. This contradicts the uniform bound
on path norms. The forest supplies arbitrarily long finite paths;
the infinite tree also permits an argument using the supremum of all
path norms. Sections~\ref{sec:coordinate-variants} and
\ref{sec:infinite-obstruction} prove these assertions for the two
completions separately.


\subsection{The other quantities measured by the methods}

A second proof method works directly with primal tests. Clipping their
cumulative sums along each exterior path produces two perturbations
with a quadratic norm cost. Testing the two endpoints on these
perturbations controls their average. Corollary~\ref{clip:paired-modulus}
gives the resulting uniform gain for perturbations in a suitable
finite-codimensional subspace. The construction and its root-shift
variants occupy Sections~\ref{clip:section}--\ref{sec:primal-additions}.

The later atomic arguments compare what survives when one changes the
normalization or retains more of the original representation. A common
head test can control several endpoint distributions with unequal
weights or a nonzero center tail. A bound on their mixture does not
automatically give the same bound on each conditional law.
Proposition~\ref{prop:endpoint}, for example, records separate
coefficient errors and boundary masses. Some alternative calculations
give weaker geometric constants: the radial bound
$(\sqrt2+\sqrt6)\sqrt{R(R-\|x\|)}$ of Proposition~\ref{thm:tail}
estimates the two errors separately, while the original-product route
incurs a fourth-root loss. These calculations record the losses of
their respective methods. Sign selection and
maximal-gain trimming instead isolate which crossing contributions
can be retained with a common sign.

These estimates have metric consequences. The companion
\emph{Distortion of countably branching diamonds from midpoint and tree energies}
\cite[Theorem~1.1]{DiamondEnergies} proves that Theorem~\ref{thm:main-tail} forces
$D^2\ge1+k/4$ for embeddings of the depth-$k$ countably branching
normalized diamond. That graph starts from an edge of length one and
replaces each edge by countably many two-edge paths of the same total
length; distortion means a common ratio between upper and lower
all-pairs Lipschitz factors. The companion proves all infinite-family
selections and graph recurrences, and also preserves the independent
radial-deficit route and direct energy methods.

The broader question whether asymptotic midpoint uniform convexity
implies an equivalent asymptotically uniformly convex norm already has a
negative answer. Baudier proved midpoint convexity for the Kadets--Werner
modification of the Bourgain--Rosenthal space, which has no equivalent
AUC norm \cite[Theorem~1 and Corollary~1]{Baudier2026}. Here the
finite-height forest is reflexive, and the infinite-height coordinate
predual is a different completion. Their exact identification statements
and local estimates are part of the proofs below. The numerical bounds
are not claims of optimality or priority.


\paragraph{Reading guide.}
Sections~\ref{sec:forest}--\ref{mol:sec:joint-budget} establish the
coordinate foundations and the main lens theorem.
Section~\ref{sec:consequences} derives the midpoint modulus; the
bounded-path arguments in Sections~\ref{sec:coordinate-variants}--
\ref{sec:infinite-obstruction} complete the principal geometric
conclusions. Sections~\ref{clip:section}--\ref{sec:primal-additions}
then develop the independent primal method, from cumulative clipping
to root shifts and bounded separated families.

The remaining sections study atomic representations in more detail.
Sections~\ref{sec:radial}--\ref{sec:joint-consequences} retain radial
and conditional estimates. Sections~\ref{presentation:constructions}--
\ref{stab:toolkit} give the gauge and polar descriptions and the common
finite-test estimates. Sections~\ref{stab:section}--\ref{stab:trim:section}
use them to compare weighted endpoint laws, original coefficients and
sign choices, with the model specified in each statement.
Sections~\ref{sec:sequential-forest}--\ref{sec:qualitative-infinite}
finish with qualitative arguments and their distinct normalization
procedures.
